\subsection{Nonorientable and real surface checks}

\begin{result}{C.9}{Known, with derivation proved}{The nonorientable comparison and real-equivariant representatives}
For a compact smooth nonorientable surface, the inclusion of
diffeomorphisms into homeomorphisms induces an isomorphism of
mapping-class groups, also with individually fixed marked points
and boundary fixed pointwise. For a closed nonorientable surface
$N_h$ with $h$ crosscaps, the unmarked identity components in both
categories have homotopy types
\[
 SO(3)\quad(h=1),\qquad S^1\quad(h=2),\qquad *\quad(h\geq3).
\]
With one fixed point $N_1$ has identity component of weak homotopy
type $S^1$; with at least two it is weakly contractible.
For $N_2$ with at least one fixed point, and for $N_h$ with
$h\geq3$ with any number of fixed points, it is weakly contractible.

Let $\tau$ be a specified smooth orientation-reversing involution
of an oriented compact surface, and retain a specified finite
$\tau$-invariant marking. The inclusion of the smooth centralizer
of $\tau$ into its topological centralizer induces an isomorphism
on component groups, with the same marking conditions. Quotient
and lift along $S\to S/\langle\tau\rangle$ give the comparison.
\end{result}
\begin{proof}
The relative handle-smoothing construction used in C.5 also works
without an orientation. Its cut arcs are straightened relative to
their endpoints, after which each complementary disk is smoothed.
The reverse implication uses the relative arc-isotopy theorem and
the disk isotopy exactly as before. This proves the component-group
assertion and its compatibility with cutting and forgetting marks.
The boundary-fixed identity-component theorem, including the two
closed nonorientable exceptions, is Theorem B.1 of
\cite{HatcherMW2018}; the corresponding topological calculation is
\cite[Theorem 1.5.2]{Hamstrom1974}.

For $N_1=\mathbb{RP}^2$, evaluate the rotation model $SO(3)$ at
an unoriented line. The base is $\mathbb{RP}^2$ and the identity
component of its stabilizer is $SO(2)$. After fixing that first
point, the punctured surface is an open M\"obius band. Evaluation
of its stabilizer circle at a second point winds twice around the
core, hence induces the injective map $\Z\xrightarrow{\,2\,}\Z$.
The evaluation exact sequence gives zero positive homotopy groups
for the identity component fixing both points. For the Klein
bottle, the circle model evaluates to the central infinite cyclic
subgroup of its fundamental group. In the usual presentation
$\langle a,b\mid aba^{-1}=b^{-1}\rangle$ this is
$\langle a^2\rangle$; its inclusion is injective. The same exact
sequence gives weak contractibility after fixing one point.
All further point additions have aspherical punctured-surface
bases, so the induction of C.5 proves the remaining assertions.

For the real-equivariant assertion, the quotient of $S$ by $\tau$
is a surface, with boundary the image of its fixed curves. Its
interior double cover records the local orientation, and the double
along the boundary reconstructs $S$; in the fixed-point-free case
this is the orientation double cover throughout. A quotient
homeomorphism respecting the retained markings has a unique
orientation-preserving lift, because its two lifts differ by the
orientation-reversing $\tau$. Along a boundary collar it is
isotopic, relative to the marked boundary points, to a collar
product. Apply the relative smoothing construction on the quotient
with that product fixed. The product lifts smoothly across the
double, so the lifted smooth representative commutes with $\tau$.
For an isotopy, make the same collar adjustment and use the
relative arc-isotopy argument on the quotient; its unique
orientation-preserving lift is an equivariant smooth isotopy.
This proves both surjectivity and injectivity on components. An
orientation-reversing centralizer element differs from an
orientation-preserving one by $\tau$, giving the full statement.

These comparisons specify the real and nonorientable groups used
for representatives. The normal lattices in A.5 arise from complex
node parameters on the oriented double: conjugation sends their
positive generators to their inverses as in A.12. No lattice for
a one-sided curve in the quotient surface is substituted for
those complex-node meridians.
\end{proof}
